% ============================================================
% Raw related-work material from chapters/agency.tex
% ("Price recursion as the rational theory of reward"), for the
% Part I literature review (chapters/litrev-agency.tex, \label{ch:litrev-agency}).
%
% The source articles have no related-work section. The only passage that is
% purely a prior-art survey is the opening of the main article
% ("Price recursion is the rational theory of reward"), reproduced verbatim
% below. In chapters/agency.tex it is commented out with
% "% [dup of thesis intro]", since Chapter 1 (intro:sec:agency) already
% surveys the same program-market lineage with the same citations.
%
% No other passage in the chapter is worth moving: the remaining citations
% (Harper 2009 and Baez 2012 on replicator dynamics as Bayesian inference;
% Liu on equilibrium thermoeconomics; the author's own LMSR-subsidy and
% market-RL papers) occur as one-clause asides inside the argument.
% ============================================================

% Source: marketgrad-theory/Price recursion is the rational theory of reward.md, opening paragraph.
Several
papers discuss RL algorithms based on sequentially auctioning a unit ``right to act'' among
agents---most recently Caspar Oesterheld's BRIA \citep{oesterheldTheoryBoundedInductive2023}
but also Baum's Hayek
Machine \citep{baumModelIntelligenceEconomy1999}, its predecessors the classifier systems
\citep{hollandPropertiesBucketBrigade1985a, schmidhuberEvolutionaryPrinciplesSelfreferential1987, schmidhuberLocalLearningAlgorithm1989} and
its successors Hayek-for-POMDPs \citep{kweeMarketBasedReinforcementLearning2001} and
Hayek with Vickrey auctions \citep{changDecentralizedReinforcementLearning2020}. In
``Market-based architectures in RL and beyond'' \citep{pallavisudhirMarketbasedArchitectures2025}
I suggested that a better market model for RL might come out of factoring the ``right to
act'' into ``goods'' axes, just like in real markets.
